\section{Experimental Comparison}

All three monitoring paradigms are implemented in a common framework sharing a uniform \textsc{Monitor} interface, which ensures that the experimental harness is paradigm-agnostic: timing experiments iterate over a list of monitor instances without any paradigm-specific branching. The shared compilation pipeline uses \texttt{ltlf2dfa} \citep{fuggitti-ltlf2dfa} to translate each \ltlf formula into a minimal DFA once at construction time.

We evaluate the monitoring paradigms along three axes, using synthetic traces throughout Experiments 1–3 to isolate per-step monitoring cost from data-dependent effects such as early-termination frequency. This follows the methodology of experiments in \citet{perotti2014neural}. 
%Here we compare their approach with the symbolic DFA and DeepDFA baselines, which were absent from the original evaluation.



\begin{comment}
All monitors share a common Python harness. Timing is measured as total wall time divided by the number of trace cells processed, following Perotti et al. (2014). Each experiment reports mean and standard deviation over multiple repetitions.
Experiment 1 — Trace length scaling. We fix the formula G(a → Fb) and vary trace length from 1,000 to 10,000 cells. This formula has neither a trap state nor an accepting sink, so early termination never fires and the per-cell cost is measured in isolation. We expect all three paradigms to show roughly flat per-cell cost as trace length grows, since each paradigm processes cells independently with no memory of past cells beyond its current state.
Experiment 2 — Formula complexity. We use the IJCNN 2014 scalability formula family ◇∨ᵢ(a₀ ∧ aᵢ) with n = 2, 4, 8, 16, 32 atoms, directly reproducing and extending the original benchmark. We expect the symbolic DFA to remain flat (the DFA size grows but per-step cost stays O(1)), RuleRunner to grow linearly in formula size (rule count is linear in parse-tree size), and DeepDFA to remain flat in per-step cost with a larger GPU-initialization overhead.
Experiment 3 — Batch size scaling. We fix the formula at n=8 atoms and vary the number of traces monitored simultaneously from 1 to 1,024. This is where DeepDFA's native GPU batching is expected to provide the clearest advantage: the symbolic DFA and RuleRunner process traces sequentially (their batch_run falls back to sequential calls), while DeepDFA processes all traces in parallel via a single batched matrix operation.
Note on RuleRunner's correctness scope: the rule-based encoding has a known structural limitation on formulas with nested temporal operators (e.g. G(a→Fb)). Experiment 1 uses such a formula for timing only — since early termination never fires, the per-cell cost measurement is valid regardless of verdict correctness. Experiments 2 and 3 use the flat-temporal IJCNN family where the encoding is verified correct. We document this limitation as a finding in the theoretical comparison (Section 4) and as a concrete argument for the automata-based paradigm's generality.
\end{comment}